\documentclass[10pt,a4paper]{amsart}
\usepackage[margin=19mm]{geometry}
\usepackage{amsmath,amssymb,mathtools}
\usepackage[scr=rsfs,cal=euler]{mathalfa}
\usepackage{xfrac}
\usepackage[unicode,hidelinks,bookmarks=false]{hyperref}
\newcommand{\ie}{i.e.\ }
\newcommand{\cf}{cf.\ }
\newcommand{\resp}{respectively\ }
\newcommand{\Subsection}{\subsection}
\providecommand{\nobreakdash}{\nobreak\hbox{-}}
\setlength{\emergencystretch}{2em}
\multlinegap0pt
\usepackage{siunitx}
\usepackage{siunitx}
\usepackage{siunitx}
\usepackage{siunitx}
\usepackage{siunitx}
\usepackage{xcolor}
\newtheorem{lemma}{Lemma}
\newtheorem{theorem}{Theorem}
\title{On the Bounds of Stopping and Cycle Numbers in the Collatz Process}
\author{Daohang Sha}
\date{Source: arXiv:2112.12962v16 (CC BY 4.0)}
\begin{document}
\maketitle

\noindent\emph{Source and licence.} On the Bounds of Stopping and Cycle Numbers in the Collatz Process. Version arXiv:2112.12962v16 (CC BY 4.0), distributed by arXiv under the Creative Commons Attribution 4.0 International licence, http://creativecommons.org/licenses/by/4.0/, as recorded in the arXiv licence metadata for this exact version and shown on the abstract page https://arxiv.org/abs/2112.12962v16. Full licence notice: \texttt{licenses/arxiv-2112.12962-CC-BY-4.0.txt}.\par\smallskip

\noindent\textit{Editorial scope.} Counted unit: the article's theorem thm:exhaustive (source0.tex lines 558--566). The article's complete original proof of it is delivered with it (source0.tex lines 568--621, one \texttt{proof} environment of 54 lines). No mathematics is written, repaired or re-derived: the statement and the proof are the article's own lines, in the article's own notation, and no retained line is substituted or re-flowed.  Retained uncounted as the source-local context the proof needs: lines 323--328; lines 419--427.  One declared adaptation: exactly 1 ancillary strings inside the retained spans are omitted (image-inclusion commands and, where the source repeats a label, the repeated label definitions) and no other line differs from the source; the figure environments, with their placement, captions and labels, are kept, and the package records the omitted strings in fidelity receipts bound to the affected bodies.  No retained line contains a citation, so the wrapper carries no bibliography and no bibliography entry.  The retained lines contain the article's own diagram code, which the wrapper typesets with the standard tikz packages; no image file is delivered and the package declares no asset dependency.  Editorial formatting only, in addition: an AMS \texttt{amsart} preamble that restates the article's own declarations and macros used by the retained lines, namely the environments \texttt{lemma}, \texttt{theorem} on separate counters, together with the standard packages those lines need.  The retained environments print in this document as Lemma 1, Theorem 1 in order of appearance, whereas the article prints the same items with the numbering of its own full document.  The complete native source of the article as distributed in the arXiv e-print for this exact version is bundled unchanged as \texttt{sources/120084/source0.tex}.
	\begin{figure}[ht]
		\centering
		
		\caption{Rooted sequence tree}
		\label{fig:seqtree}
	\end{figure}

\begin{lemma}\label{lem:3mod4}

Let $m \equiv 3 \pmod{4}$. Then $m$ admits a unique representation of the form
$$ m = 3 + \sum_{n_i\ge2} 2^{n_i}, \qquad n_i \ge 2, $$
where the exponents $n_i$ are distinct. Conversely, every finite sum of the form
$$ 3 + \sum_{n_i\ge2} 2^{n_i}, \qquad n_i \ge 2,$$
yields an integer $m \equiv 3 \pmod{4}$.

\end{lemma}

\begin{theorem}[Exhaustivity for $m \equiv 3 \pmod{4}$]\label{thm:exhaustive}
	The construction in Figure~\ref{fig:seqtree} generates each integer
	$m \equiv 3 \pmod{4}$ exactly once. Equivalently, the correspondence
	\[
	q \;\longleftrightarrow\; m
	\]
	defined by Figure~\ref{fig:seqtree} is a bijection between admissible sequences $q$
	and integers $m \equiv 3 \pmod{4}$.
\end{theorem}

\begin{proof}
	We prove the statement by showing that the constructions
	$q \mapsto m$ and $m \mapsto q$ are well-defined and inverse to each other.
	
	\medskip
	\noindent\textbf{Forward direction ($q \mapsto m$).}
	Starting from the root $q = 11$, we have
	\[
	F_{11}(4i - 1) = 3^{2} i - 1.
	\]
	Any admissible sequence $q = 11\cdots$ corresponds to a path in the rooted
	sequence tree shown in Figure~\ref{fig:seqtree}. Along this path, the Collatz
	process, $F_q(4i - 1)=3^ry+F_q(m)$, yields successive decompositions
	\[
	i = 2j \text{ or } 2j + 1,\quad
	j = 2k \text{ or } 2k + 1,\quad
	k = 2l \text{ or } 2l + 1,\ \ldots,\ 
	x = 2y \text{ or } 2y + 1.
	\]
	Consequently, we may write
	\[
	4i - 1 = 2^{s} y + \sum_{n_i \ge 2} 2^{n_i} + 3
	= 2^{s} y + m,
	\]
	which implies
	\[
	m \equiv 3 \pmod{4}.
	\]
	
	\medskip
	\noindent\textbf{Backward direction ($m \mapsto q$).}
	Let $m$ be any integer satisfying $m \equiv 3 \pmod{4}$. By
	Lemma~\ref{lem:3mod4}, $m$ admits a representation of the form
	\[
	m = \sum_{n_i \ge 2} 2^{n_i} + 3.
	\]
	Writing $4i - 1 = 2^ry + m$ and using the powers $2^{n_i}$ in the above decomposition,
	we recover successive relations
	\[
	i = 2j \text{ or } 2j + 1,\quad
	j = 2k \text{ or } 2k + 1,\quad
	k = 2l \text{ or } 2l + 1,\ \ldots,
	\]
	which uniquely determine an admissible sequence
	\[
	q = 11\cdots
	\]
	can be obtained via the inverse Collatz process and the mapping $F_q(4i - 1)=3^ry+F_q(m)$.
	
	\medskip
	Since the mappings $q \mapsto m$ and $m \mapsto q$ are inverses of each other,
	the correspondence between admissible sequences $q$ and integers
	$m \equiv 3 \pmod{4}$ is bijective.
\end{proof}

\end{document}
